\section{System Overview}
\label{sec:overview}

\subsection{Problem Formulation}
\label{sec:problem}

\textbf{Estimation problem.}
We estimate the continuous-time 6-DoF body trajectory (world-frame
position $\vp_w(t)\in\mathbb{R}^3$ and orientation $\vR(t)\in\SO$)
together with the constant \ac{imu} biases $\vb_a,\vb_g$, from asynchronous
measurements over the estimation interval (the full trajectory in batch
mode, a sliding window online): radar Doppler velocities ($\sim$11\,Hz),
accelerometer and gyroscope samples (1\,kHz), and a low-rate heading
prior.  We adopt a \emph{continuous-time} representation because the
sensors are asynchronous and run at disparate rates: a trajectory defined
for every $t$ is evaluated at each measurement's exact timestamp, and its
derivatives (velocity $\dot{\vp}_w$, acceleration $\ddot{\vp}_w$, and
body angular velocity $\vom_b$) follow in closed form, so no
inter-sample interpolation or rigid-keyframe approximation is required.

\textbf{Parameterization.}
The trajectory is a B-spline (\autoref{sec:state}): the position is a
quintic uniform B-spline with control points
$\{\mathbf{P}_i\}\subset\mathbb{R}^3$, and the orientation a cumulative
$\SO$ B-spline with \emph{absolute} rotation knots
$\{\vR_i\}\subset\SO$ (basalt convention~\cite{usenko2020basalt}).
This renders the otherwise infinite-dimensional trajectory
finite-dimensional; the control points and the \ac{imu} biases form the
estimated state
\begin{equation}
  \mathcal{X} = \bigl\{\,\mathbf{P}_i,\ \vR_i,\ \vb_a,\ \vb_g\,\bigr\}.
  \label{eq:state}
\end{equation}
The batch solve additionally frees the single observable extrinsic degree
of freedom (the radar pitch); the sliding window fixes it at the measured
value (\autoref{sec:state}).

\textbf{Estimator.}
Treating the measurement noise as zero-mean and independent, the
\ac{map} estimate is the minimizer of the negative
log-posterior, a weighted nonlinear least-squares problem over the
measurement residuals, the spline regularizers, and the priors:
\begin{equation}
\begin{aligned}
  \hat{\mathcal{X}} = \operatorname*{arg\,min}_{\mathcal{X}}
    & \sum_{k}\rho_\delta\!\bigl(r_{D,k}^2\bigr)
    + \lambda_a\!\sum_{m}\lVert r_{a,m}\rVert^2
    + \lambda_g\!\sum_{m}\lVert r_{g,m}\rVert^2 \\
    & + \lambda_\psi\!\sum_{n} r_{\psi,n}^2
    + E_{\text{reg}}(\mathcal{X})
    + E_{\text{prior}}(\mathcal{X}),
\end{aligned}
\label{eq:map}
\end{equation}
where $r_D,r_a,r_g,r_\psi$ are the radar-Doppler, accelerometer,
gyroscope, and heading residuals of \autoref{sec:models};
$\rho_\delta$ is the Huber kernel ($\delta{=}1.0$\,m/s) that bounds the
influence of heavy-tailed Doppler outliers; the $\lambda_\bullet$ are the
(rate-adaptive) information weights of \autoref{sec:models};
$E_{\text{reg}}$ collects the position min-snap and orientation
smoothness regularizers (\autoref{sec:reg}); and $E_{\text{prior}}$
holds the boundary and, online, the marginalization priors
(\autoref{sec:sw}).  We minimize~\eqref{eq:map} with
Levenberg--Marquardt (Ceres) using analytic Jacobians.  The fixed-lag
estimator solves it over a 3\,s window and, on each stride, marginalizes
the states leaving the window into $E_{\text{prior}}$ via the Schur
complement (\autoref{sec:sw}), keeping the per-window cost bounded.

\subsection{Hardware Platform}
The system runs on a quadrotor with a TI IWR6843AOPEVM
mmWave radar (60--64\,GHz \ac{fmcw}), a Pixhawk flight
controller \ac{imu}, and a Vicon \ac{mocap} system used for ground
truth evaluation.  The radar is mounted upside-down at a nominal $30\degree$ downtilt by
\ac{cad}; the as-built angle is a few degrees lower (\autoref{sec:results}),
and the extrinsic rotation is $[180\degree, 27.5\degree, 0\degree]$
(roll, pitch, yaw), translation $[0.08, 0.02, -0.01]$\,m in the body
frame.  The radar boresight lies along the body forward axis, so the
Euler pitch is numerically the physical boresight downtilt.  The radar operates at $\sim$11\,Hz with 0.049\,m/s Doppler
resolution (best-velocity firmware configuration).  Raw \ac{imu} rate
is $\sim$993\,Hz; both modalities are used at full rate.

\subsection{Pipeline}
\autoref{fig:pipeline} shows the pipeline.  A three-stage sensor-only
initialization (P1--P3, detailed in \autoref{sec:init}) builds the
initial trajectory the nonlinear solve needs, from radar and \ac{imu}
data alone; the same raw measurements then also enter the solver directly
as factors (\autoref{fig:pipeline}, dashed).  \ac{mocap} supplies only the
yaw heading prior (a magnetometer surrogate, \autoref{sec:models}) and the
final \ac{rmse} evaluation; P1--P3 and the position/attitude estimate use
no external pose reference.

\begin{figure}[t]
\centering
\scalebox{0.78}{%
\begin{tikzpicture}[
  font=\small,
  box/.style={rectangle, rounded corners=3pt, draw, minimum width=1.55cm,
              minimum height=0.52cm, align=center},
  sensor/.style={box, fill=blue!12},
  stage/.style={box, fill=orange!20},
  solver/.style={box, fill=green!15, minimum width=1.6cm, minimum height=0.95cm},
  arr/.style={-{Latex[length=2mm]}, thick},
  darr/.style={-{Latex[length=2mm]}, thick, dashed, gray!70},
  ]
  %% Nodes: sensor column, two init stages, merge/P3 stage, solver, output
  \node[sensor] (imu)    at (0.0,  1.1) {\ac{imu}\\1\,kHz};
  \node[sensor] (radar)  at (0.0, -0.5) {Radar\\$\sim$11\,Hz};
  \node[stage]  (p1)     at (2.2,  2.0) {P1: Gyro\\${\to}R(t)$};
  \node[stage]  (p2)     at (2.2,  0.2) {P2: \ac{wls}\\${\to}\mathbf{p}(t)$};
  \node[stage]  (p3)     at (4.1,  1.1) {P3: Priors\\${\to}$ init};
  \node[solver] (solver) at (5.9,  1.1) {C++ Ceres\\(batch/SW)};
  \node[box, fill=gray!10] (out) at (7.7, 1.1)
        {$\hat{\mathbf{p}},\hat{R},\hat{\mathbf{v}}$};

  %% Init pipeline arrows (solid, no crossings by construction)
  % IMU → P1 (upper-right diagonal)
  \draw[arr] (imu.north east) -- (p1.west);
  % IMU → P2 (lower-right diagonal; meets Radar→P2 at P2, no earlier crossing)
  \draw[arr] (imu.south east) -- (p2.west);
  % Radar → P2 (upper-right; both end at P2 → no crossing)
  \draw[arr] (radar.east)     -- (p2.west);
  % P1 and P2 → P3 (converge; meet at P3, no earlier crossing)
  \draw[arr] (p1.east)        -- (p3.north west);
  \draw[arr] (p2.east)        -- (p3.south west);
  % P3 → Solver
  \draw[arr] (p3.east)        -- (solver.west);
  % Solver → Output
  \draw[arr] (solver.east)    -- (out.west);

  %% Raw-data feeds (dashed, routed above/below init block to avoid crossings)
  \draw[darr] (imu.east)   to[out=10,  in=170] (solver.north west);
  \draw[darr] (radar.east) to[out=350, in=190] (solver.south west);
\end{tikzpicture}%
}% end scalebox
\caption{System pipeline. Solid: initialization flow (P1--P3 seed
the optimizer). Dashed: raw sensor data fed as measurement factors.
Radar returns first pass a \ac{ransac} ego-velocity prefilter
(\autoref{sec:models}) that rejects elevation-biased outliers before
both P2 and the solver. \ac{mocap} supplies only the yaw heading prior
(magnetometer surrogate) and \ac{rmse} evaluation.}
\label{fig:pipeline}
\end{figure}
